%% =================================================================
%% Wei, Mei, Zhao, Xiao, Sun, Zhang, Guo.
%% "Nondestructively identifying the mechanical behavior of soft
%%  tissues using surface deformation with an explicit inverse
%%  approach."
%% Appl. Math. Modelling 134 (2024) 126--147.
%%
%% Reconstruction of page 13 (printed p. 138) as study notes.
%% Generated by: reproduce-multiple-pages-basic-tex (batch 13-16)
%%
%% Structure: mirrors pages 12 — two full-page result figures
%% (Fig. 14 for Sample 1 and Fig. 15 for Sample 2, both 5% noise),
%% no prose, no equations, no tables.
%% Math: none on this page.
%% Figures 14 & 15: placeholder boxes with captions + visual descriptions.
%% =================================================================

\documentclass[11pt]{article}

\usepackage[margin=1in]{geometry}
\usepackage{amsmath, amssymb}
\usepackage{mathrsfs}
\usepackage{bm}
\usepackage{xcolor}
\usepackage{fancyhdr}

\pagestyle{fancy}
\fancyhf{}
\lhead{\small Y.~Mei et al.}
\rhead{\small Applied Mathematical Modelling 134 (2024) 126--147}
\cfoot{\small 138 $\to$ \arabic{page}}
\renewcommand{\headrulewidth}{0.4pt}

\setcounter{page}{1}

\begin{document}

% ---------- Figure 14 placeholder ----------
\noindent
\begin{center}
\fbox{\parbox{0.92\textwidth}{\centering\vspace{1em}
\textbf{[Figure 14 --- placeholder]}\\[0.8em]
Six-panel side-by-side comparison of the \textbf{Explicit Inverse
Method} (left column) vs.\ \textbf{Implicit Inverse Method} (right
column) on \textbf{Sample 1} (the homogeneous target, true
SM $= 3.00$\,MPa), with \textbf{5\,\% noise} added to the measured
surface displacements, using 3, 6, and 9 displacement datasets.\\[0.4em]
\textbf{Explicit (left column)}: the higher noise level corrupts
the homogeneous recovery. Colorbar ranges:
(a) 3 fields, SM $\in [2.9641, 3.0096]$\,MPa --- visible spurious
\textbf{bilayer artifact} with a blue ``valley'' cutting through
an otherwise red field (the homogeneous target should be uniform);
(b) 6 fields, SM $\in [3.0014, 3.0099]$\,MPa --- wavy spurious
bilayer pattern, red top over blue bottom;
(c) 9 fields, SM $\in [3.0001, 3.0001]$\,MPa --- uniform blue, near
perfect. The 9-field case recovers cleanly while 3-field and
6-field hallucinate structure that isn't there.\\[0.3em]
\textbf{Implicit (right column)}:
(d) 3 fields, SM $\in [2.7156, 3.0220]$\,MPa --- mostly uniform red
with a small localized blue blob;
(e) 6 fields, SM $\in [3.0027, 3.0060]$\,MPa --- visible vertical
gradient/striping;
(f) 9 fields, SM $\in [2.9999, 3.0002]$\,MPa --- almost uniform with
faint artifact pattern.\\[0.3em]
\textbf{Takeaway}: under 5\,\% noise on a homogeneous target, the
explicit method's 3-field and 6-field reconstructions show
\emph{spurious bilayer patterns} that don't exist in the ground
truth. This is a failure mode worth flagging: the method can
invent structure under noise when very few displacement datasets
are available. Nine fields restores clean recovery.
\vspace{1em}}}
\end{center}
\vspace{0.3em}
\begin{center}
\textbf{Fig. 14.} The reconstructed results with varying 3, 6 and
9 surface displacement datasets corresponding to Sample 1 (5\,\%
noise is introduced into the measured data, Unit: MPa).
\end{center}

\vspace{0.8em}

% ---------- Figure 15 placeholder ----------
\noindent
\begin{center}
\fbox{\parbox{0.92\textwidth}{\centering\vspace{1em}
\textbf{[Figure 15 --- placeholder]}\\[0.8em]
Six-panel side-by-side comparison of the \textbf{Explicit Inverse
Method} (left column) vs.\ \textbf{Implicit Inverse Method} (right
column) on \textbf{Sample 2} (the two-layer target: background
$\mu^0 \approx 3$\,MPa, top layer $\mu^1 \approx 10$\,MPa),
with \textbf{5\,\% noise} added, using 3, 6, and 9 displacement
datasets.\\[0.4em]
\textbf{Explicit (left column)}: bilayer structure visible in all
three panels. Colorbar ranges:
(a) 3 fields, SM $\in [1.00, 8.91]$\,MPa --- notably, the colorbar
max is \textbf{11\,\% below the nominal 10\,MPa} top layer. This
is an \emph{undershoot}, opposite to the usual overshoot failure
mode. The top layer is not fully recovered with only 3 fields
under 5\,\% noise;
(b) 6 fields, SM $\in [1.00, 10.87]$\,MPa --- slight overshoot,
wavy interface;
(c) 9 fields, SM $\in [1.00, 10.16]$\,MPa --- near-nominal with
slight overshoot.\\[0.3em]
\textbf{Implicit (right column)}:
(d) 3 fields, SM $\in [1.00, 14.20]$\,MPa --- \textbf{42\,\%
overshoot} above the 10\,MPa nominal, visible with significant
green/yellow patches intruding into the top layer;
(e) 6 fields, SM $\in [1.00, 11.42]$\,MPa;
(f) 9 fields, SM $\in [1.00, 11.43]$\,MPa --- surprising that 9
fields is not clearly better than 6 fields for the implicit
method on this sample.\\[0.3em]
\textbf{Takeaway}: the 3-field explicit case \emph{undershoots}
while the 3-field implicit case \emph{overshoots} by 42\,\% ---
the two methods fail in opposite directions under low-data +
5\,\% noise. More displacement fields close the gap for explicit
but not cleanly for implicit.
\vspace{1em}}}
\end{center}
\vspace{0.3em}
\begin{center}
\textbf{Fig. 15.} The reconstructed results with varying 3, 6 and
9 surface displacement datasets corresponding to Sample 2 (5\,\%
noise is introduced into the measured data, Unit: MPa).
\end{center}

% [page ends here — continues on page 14]

\end{document}
